% Abstract em inglês
Mercury speciation is an analytical technique used to identify and quantify the
different chemical species of this element in environmental samples, such as
elemental mercury (Hg$^0$) and methylmercury. The sample preparation procedure
requires strict control of a sequence of chemical steps --- derivatization with
sodium tetraethylborate, cryofocusing in liquid nitrogen at approximately
$-196\,\celsius$, controlled heating of the U-tube up to $230\,\celsius$, and
atomization at $700\,\celsius$ --- so that the quality of the temperature ramp
applied to the U-tube directly determines the correct separation and
identification of the species.

This work presents the design and implementation of an automation and
supervision system for sample preparation in mercury speciation, conceived to
replace a legacy LabVIEW-based system. The system adopts a master-slave
architecture composed of a Raspberry Pi running the high-level logic in Python
(finite state machine, proportional-integral-derivative --- PID --- control and
parameter persistence), an Arduino Uno operating as a real-time data
acquisition module responsible for driving valves and furnaces, reading
thermocouples through the SPI interface and running a safety watchdog, and a
Web human-machine interface (HMI) developed in React and TypeScript, with
real-time monitoring over WebSocket, a timing diagram, trend charts and
automatic and manual modes.

The process was modeled as a finite state machine with four actuation phases
($T_0$ to $T_3$) and an actuator matrix (solenoid valves SV1 to SV5,
peristaltic pump and furnaces). Temperature control employs a mixed PID
strategy: below $0\,\celsius$, the manipulated variable is computed as the ratio
between the operator-required heating rate and the system rate, whereas above
this threshold a closed-loop PID controller tracks a dynamic ramp setpoint up to
$230\,\celsius$; Furnace 2 is regulated at a fixed setpoint of $700\,\celsius$.
Communication between the Raspberry Pi and the Arduino is carried out through
JSON packets over a serial link at 4~Hz, with a parametrization handshake and a
double watchdog that forces the return to the safe state in case of
communication loss.

The development followed a test-driven approach, validated by unit,
integration, functional and safety tests executed in a simulation environment
without physical hardware. The results show that the proposed architecture
meets the defined functional and safety requirements, ensuring repeatability of
the control logic, auditable persistence of the method parameters and automatic
safe shutdown in the presence of communication failures. Bench validation, with
the calibration of the U-tube heating curve and the experimental verification of
the temperature ramp, is presented as a follow-up stage of this work.

\vspace{0.5cm}
\noindent\textbf{Keywords:} automation; PID control; mercury speciation;
finite state machine; human-machine interface; embedded system.
